\chapter*{Préface}
\addcontentsline{toc}{chapter}{Préface}
\markboth{Préface}{Préface}

\section*{Interaction complexe entre deux machines simples}
\phantomsection
\addcontentsline{toc}{section}{Interaction complexe entre deux machines simples}

Ce document rassemble le vaste matériel technique — pour de futures publications — des résultats obtenus et des développements mathématico-physiques réalisés entre les années 2023 et 2026 au moyen d'outils de Large Language Models (LLM). De ce processus émergent ce que nous appelons, respectivement, holographie modulaire triadique et mécanique dimensionnelle (HMT–MD). Cependant, l'origine de ce projet remonte au début du XXIe siècle, plus précisément à l'étude de la découverte sous-jacente à l'artefact inventé par Oumar Haidara Fall en 2001. Il convient de préciser ici que cet artefact a été porté à la connaissance des diverses autorités académiques et civiles compétentes dans les lieux d'origine respectifs des deux auteurs.

Le matériel accumulé et les connaissances acquises progressivement au cours de plus de deux décennies d'observation, d'étude et d'analyse, tant de l'artefact que de sa phénoménologie, ont constitué la base opératoire et informationnelle que nous avons ensuite appliquée aux modèles d'intelligence artificielle.

Ainsi, le résultat HMT–MD que nous présentons ici, avec ses dérivations et ses démonstrations, est précédé d'une généalogie empirique qui, \textit{a priori}, pourrait sembler, dans le meilleur des cas, contre-intuitive au technicien spécialiste — puisqu'elle contrevient aux postulats et aux prédictions du cadre théorico-interprétatif dans lequel, par tradition historique, s'inscrit l'artefact, à savoir la mécanique classique —. Toutefois, sa vérification métrologique s'avère décisive. En outre, en surmontant les préjugés cognitifs ou les biais de confirmation, elle devient directement démontrable par les faits.

Enfin, il conviendrait de considérer ce qui ne cesse de surprendre comme paradoxe historique : alors que nous commencions à raisonner avec des « machines pensantes » à « l'ère de l'intelligence artificielle », nous avons aussi eu la sérendipité de dévoiler une interaction inédite dans la mécanique ordinaire et nous nous sommes ainsi trouvés contraints de réexaminer la relation entre deux machines simples : le levier et la roue.
